\begin{proofbox}
\textbf{\textcolor{CProof}{证明目标}}
对 $C:y^2=2px\ (p>0)$，$P(x_0,y_0)$ 满足 $y_0^2-2px_0>0$，证明过 $P$ 所引两切线的切点弦为
\[\boxed{y_0y=p(x+x_0)}.\]
以下四法全部只用高中数学中的\textbf{点斜式、切线方程、方程相加与相减、二次方程和判别式}；第四法不把切点弦公式当作已知结论。

\medskip

\par\noindent\begin{minipage}{\linewidth}
\textbf{\textcolor{CProof}{方法一｜门槛最低：相减求斜率、相加找中点关系、点斜式还原（学生自推法）}}
\par\smallskip\begin{center}\begin{tikzpicture}[x=1cm,y=1cm,>=Stealth,font=\small]
\path[use as bounding box] (-2,-1.8) rectangle (10,6.4);
\node[anchor=west,font=\small\bfseries] at (-1.8,6) {斜弦的中点仍与 $P$ 同高，而 $P$ 不在弦上};
\draw[->,gray] (-1.6,0)--(6.6,0) node[right] {$x$};
\draw[->,gray] (0,-1.3)--(0,5.1) node[right] {$y$};
\draw[blue,thick,domain=-1.2:5.1,samples=81] plot ({\x*\x/4},\x);
\draw[orange!85!black,thick] (-1,2)--({3+2*sqrt(2)},{2+2*sqrt(2)});
\draw[orange!85!black,thick] (-1,2)--({3-2*sqrt(2)},{2-2*sqrt(2)});
\draw[red!75!black,very thick] ({3-2*sqrt(2)},{2-2*sqrt(2)})--({3+2*sqrt(2)},{2+2*sqrt(2)});
\draw[teal!80!black,dashed,thick] (-1.5,2)--(3.7,2);
\fill (-1,2) circle(2pt) node[above left] {$P(-1,2)$};
\fill (3,2) circle(2pt) node[below right] {$M(3,2)$};
\fill ({3+2*sqrt(2)},{2+2*sqrt(2)}) circle(2pt) node[above left] {$A$};
\fill ({3-2*sqrt(2)},{2-2*sqrt(2)}) circle(2pt) node[below left,xshift=-4pt] {$B$};
\node[anchor=north west,align=left,text width=3cm] at (7,4.8) {$y^2=4x$\\[6pt]红弦：$y=x-1$\\[10pt]$y_A=2+2\sqrt2$\\$y_B=2-2\sqrt2$\\[10pt]$y_M=\dfrac{y_A+y_B}{2}$\\[4pt]$\phantom{y_M}=2=y_0$};
\end{tikzpicture}\end{center}\par\smallskip
\end{minipage}\par
\textbf{思路：} 求直线只需一个点和一个斜率。两切线有公共点 $P$，恰好可以列出两式；\textbf{相减找方向，相加找中点的关系}，最后使用点斜式。
\begin{description}[style=nextline,leftmargin=2.8em,labelsep=0.5em,itemsep=0.5em]
\item[\textbf{第1步：列出公共点 $P$ 满足的两个等式。}]
设切点 $A(x_1,y_1),B(x_2,y_2)$，弦中点 $M(x_m,y_m)$。
因为 $A,B$ 处的切线均经过 $P$，所以
\[
 \begin{cases}
 y_0y_1=p(x_0+x_1), &(1)\\
 y_0y_2=p(x_0+x_2). &(2)
 \end{cases}
\]
\textbf{理由：} 切点 $A$ 处切线为 $y_1y=p(x+x_1)$，把 $P$ 的坐标代入即可；另一点同理。
\item[\textbf{第2步：两式相减，得到弦的斜率。}]
\[
 (1)-(2):\quad y_0(y_1-y_2)=p(x_1-x_2).
\]
若 $y_0\ne0$，又因为 $A,B$ 为不同切点，所以 $x_1\ne x_2$，于是
\[
 \boxed{k_{AB}=\frac{y_1-y_2}{x_1-x_2}=\frac p{y_0}}.
\]
\textbf{几何意义：} 两式相减消去共同的 $x_0$，得到两端点坐标之差，即\textbf{弦的方向}。
\item[\textbf{第3步：两式相加，得到中点坐标之间的关系。}]
\[
 (1)+(2):\quad y_0(y_1+y_2)=p(2x_0+x_1+x_2).
\]
由于 $x_1+x_2=2x_m$、$y_1+y_2=2y_m$，所以
\[
 \boxed{y_0y_m=p(x_0+x_m)}.
\]
\textbf{提醒：} 此处\textbf{没有分别算出 $x_m,y_m$}，而是找到了中点坐标满足的等式，这已经足够。
\item[\textbf{第4步：使用点斜式，消去中点坐标。}]
在 $y_0\ne0$ 的情况下，直线 $AB$ 经过 $M$ 且斜率已知：
\[
 \begin{aligned}
 y-y_m&=\frac p{y_0}(x-x_m)\\
 \Longrightarrow\quad y_0y-y_0y_m&=p(x-x_m)\\
 \Longrightarrow\quad y_0y-p(x_0+x_m)&=p(x-x_m)\\
 \Longrightarrow\quad &\boxed{y_0y=p(x+x_0)}.
 \end{aligned}
\]
\textbf{观察：} 中点横坐标 $x_m$ 在最后一步恰好抵消。
\item[\textbf{第5步：补上 $y_0=0$ 的竖直弦。}]
此时从 (1)、(2) 直接得到 $x_1=x_2=-x_0$，两切点在同一条竖直直线上，故
\[\boxed{AB:\ x=-x_0}\]
这也正是 $y_0y=p(x+x_0)$ 在 $y_0=0$ 时的结果。不能写 $p/y_0$。
\end{description}

\textbf{继续挖掘：能否真的求出中点 $M$？} 两切点还满足 $y_1^2=2px_1$、$y_2^2=2px_2$。两式相减：
\[
 (y_1+y_2)(y_1-y_2)=2p(x_1-x_2).
\]
代入第2步的 $p(x_1-x_2)=y_0(y_1-y_2)$，且两切点不同所以 $y_1\ne y_2$，得到
\[
 y_1+y_2=2y_0\quad\Longrightarrow\quad\boxed{y_m=y_0}.
\]
再代入第3步 $y_0y_m=p(x_0+x_m)$，得到
\[
 \boxed{M\left(\frac{y_0^2}{p}-x_0,\,y_0\right)}.
\]
\textbf{几何意义：} 弦中点与两切线交点 $P$ 的纵坐标相同。\textbf{这一步不是证明切点弦方程所必需的，而是高价值推论。}

\medskip
\par\noindent\begin{minipage}{\linewidth}
\textbf{\textcolor{CProof}{方法二｜速度最快：相减求斜率，直接使用切点的点斜式（考场速写）}}
\par\smallskip\begin{center}\begin{tikzpicture}[x=1cm,y=0.86cm,>=Stealth,font=\small]
\path[use as bounding box] (-2,-3.5) rectangle (10,4.2);
\node[anchor=west,font=\small\bfseries] at (-1.8,3.8) {方法二：把注意力从切线上点 $P$ 转向切点 $Q$};
\draw[->,gray] (-1.7,0)--(3.3,0) node[right] {$x$};
\draw[->,gray] (0,-3.1)--(0,3.1) node[above] {$y$};
\node[below left] at (0,0) {$O$};
\draw[blue,thick,domain=-3:3,samples=81] plot ({\x*\x/4},\x);
\draw[orange!85!black,thick] (-1.4,-.4)--(1.8,2.8);
\draw[orange!85!black,thick] (-1.4,.4)--(1.8,-2.8);
\draw[red!75!black,very thick] (1,-2.7)--(1,2.7);
\fill (-1,0) circle (2pt) node[above left,yshift=12pt] {$P(x_0,y_0)$};
\fill (1,2) circle (2pt) node[below right,xshift=5pt,yshift=-3pt] {$Q(u,v)$};
\fill (1,-2) circle (2pt) node[above right,xshift=5pt,yshift=3pt] {$B$};

\node[anchor=north west,align=left,text width=5.7cm] at (4,2.7) {观察上方切点 $Q$：\\[6pt]橙线过 $P$，所以\\$vy_0=p(x_0+u)$。\\[6pt]交换因子、加数次序：\\$y_0v=p(u+x_0)$。\\[6pt]这说明 $Q$ 在红线上。};
\node[anchor=west,blue] at (-1.6,-3.25) {$y^2=4x$（精确示例，$p=2$）};
\end{tikzpicture}\end{center}\par\smallskip
\end{minipage}\par
当 $y_0\ne0$，由两切点处切线都经过 $P$，相减即得
\[
 k_{AB}=\frac{y_1-y_2}{x_1-x_2}=\frac p{y_0}.
\]
直接取弦上已有的切点 $A(x_1,y_1)$，写点斜式：
\[
 \begin{aligned}
 y-y_1&=\frac p{y_0}(x-x_1)\\
 \Longrightarrow\quad y_0y-y_0y_1&=p(x-x_1)\\
 \overset{y_0y_1=p(x_0+x_1)}{\Longrightarrow}\quad
 \boxed{y_0y=p(x+x_0)}.
 \end{aligned}
\]
\textbf{为何更快：} 与方法一相比，\textbf{无需相加、无需设置中点}，斜率求出后直接使用切点 $A$。若 $y_0=0$，由两切线过 $P$ 的等式得 $x_1=x_2=-x_0$，直接写竖直弦。

\medskip
\par\noindent\begin{minipage}{\linewidth}
\textbf{\textcolor{CProof}{方法三｜步骤最少：两个切点都满足同一个一次方程（两步）}}
\par\smallskip\begin{center}\begin{tikzpicture}[x=1cm,y=0.86cm,>=Stealth,font=\small]
\path[use as bounding box] (-2,-3.5) rectangle (10,4.2);
\node[anchor=west,font=\small\bfseries] at (-1.8,3.8) {方法一：分别沿两条切线代入同一个 $P$};
\draw[->,gray] (-1.7,0)--(3.3,0) node[right] {$x$};
\draw[->,gray] (0,-3.1)--(0,3.1) node[above] {$y$};
\node[below left] at (0,0) {$O$};
\draw[blue,thick,domain=-3:3,samples=81] plot ({\x*\x/4},\x);
\draw[orange!85!black,thick] (-1.4,-.4)--(1.8,2.8);
\draw[orange!85!black,thick] (-1.4,.4)--(1.8,-2.8);
\draw[red!75!black,very thick] (1,-2.7)--(1,2.7);
\fill (-1,0) circle (2pt) node[above left,yshift=12pt] {$P$};
\fill (1,2) circle (2pt) node[below right,xshift=5pt,yshift=-3pt] {$A(x_1,y_1)$};
\fill (1,-2) circle (2pt) node[above right,xshift=5pt,yshift=3pt] {$B(x_2,y_2)$};

\node[anchor=north west,align=left,text width=5.7cm] at (4,2.7) {$P\in PA$ 给出\\$y_0y_1=p(x_1+x_0)$。\\[8pt]$P\in PB$ 给出\\$y_0y_2=p(x_2+x_0)$。\\[8pt]两个切点都在红线上。};
\node[anchor=west,blue] at (-1.6,-3.25) {$y^2=4x$（精确示例，$p=2$）};
\end{tikzpicture}\end{center}\par\smallskip
\end{minipage}\par
\textbf{第1步：} 对任意切点 $Q(u,v)$，其切线 $vy=p(x+u)$ 经过 $P$，所以
\[
 vy_0=p(x_0+u)\quad\Longleftrightarrow\quad y_0v=p(u+x_0).
\]
\textbf{第2步：} 这对不同切点 $A,B$ 都成立，因此两点都在直线 $y_0y=p(x+x_0)$ 上，故
\[\boxed{AB:\ y_0y=p(x+x_0)}.\]
\textbf{优势：} 不需计算斜率，不会遇到除以 $y_0$ 的问题；从已有切线公式出发，真正只用两步。

\textbf{补充反向检验（原有思路保留）：} 若先写候选直线 $L:y_0y=p(x+x_0)$，联立 $x=y^2/(2p)$，得
\[
 y^2-2y_0y+2px_0=0,
\]
由已知 $y_0^2-2px_0>0$ 知它有两个不同实交点。任一交点 $Q(u,v)$ 都满足 $vy_0=p(x_0+u)$，说明 $Q$ 处切线经过 $P$，因此这两个交点确实就是切点。\textbf{这是验证，不是再增加一种必背的证明法。}
\par\smallskip\begin{center}\begin{tikzpicture}[x=1cm,y=0.86cm,>=Stealth,font=\small]
\path[use as bounding box] (-2,-3.5) rectangle (10,4.2);
\node[anchor=west,font=\small\bfseries] at (-1.8,3.8) {方法三：先找红线交点，再检验两条橙色切线};
\draw[->,gray] (-1.7,0)--(3.3,0) node[right] {$x$};
\draw[->,gray] (0,-3.1)--(0,3.1) node[above] {$y$};
\node[below left] at (0,0) {$O$};
\draw[blue,thick,domain=-3:3,samples=81] plot ({\x*\x/4},\x);
\draw[orange!85!black,thick] (-1.4,-.4)--(1.8,2.8);
\draw[orange!85!black,thick] (-1.4,.4)--(1.8,-2.8);
\draw[red!75!black,very thick] (1,-2.7)--(1,2.7);
\fill (-1,0) circle (2pt) node[above left,yshift=12pt] {$P$};
\fill (1,2) circle (2pt) node[below right,xshift=5pt,yshift=-3pt] {$A$};
\fill (1,-2) circle (2pt) node[above right,xshift=5pt,yshift=3pt] {$B$};

\node[anchor=north west,align=left,text width=5.7cm] at (4,2.7) {先画候选直线\\$L:y_0y=p(x+x_0)$，\\与抛物线交于 $A,B$。\\[8pt]再在 $A,B$ 处作切线，\\验证两条切线都经过 $P$。\\[8pt]顺序：交点 $\to$ 切点。};
\node[anchor=west,blue] at (-1.6,-3.25) {$y^2=4x$（精确示例，$p=2$）};
\end{tikzpicture}\end{center}

\medskip
\par\noindent\begin{minipage}{\linewidth}
\textbf{\textcolor{CProof}{方法四｜常规通用：设直线、联立、判别式（完整代数法）}}
\par\smallskip\begin{center}\begin{tikzpicture}[x=1cm,y=0.86cm,>=Stealth,font=\small]
\path[use as bounding box] (-2,-3.5) rectangle (10,4.2);
\node[anchor=west,font=\small\bfseries] at (-1.8,3.8) {方法四：过 $P$ 的直线相切时，交点方程有重根};
\draw[->,gray] (-1.7,0)--(3.3,0) node[right] {$x$};
\draw[->,gray] (0,-3.1)--(0,3.1) node[above] {$y$};
\node[below left] at (0,0) {$O$};
\draw[blue,thick,domain=-3:3,samples=81] plot ({\x*\x/4},\x);
\draw[orange!85!black,thick] (-1.4,-.4)--(1.8,2.8);
\draw[orange!85!black,thick] (-1.4,.4)--(1.8,-2.8);
\draw[red!75!black,very thick] (1,-2.7)--(1,2.7);
\fill (-1,0) circle (2pt) node[above left,yshift=12pt] {$P$};
\fill (1,2) circle (2pt) node[below right,xshift=5pt,yshift=-3pt] {$Q\ (m=1)$};
\fill (1,-2) circle (2pt) node[above right,xshift=5pt,yshift=3pt] {$Q\ (m=-1)$};

\node[anchor=north west,align=left,text width=5.7cm] at (4,2.7) {设 $\ell:x=x_0+m(y-y_0)$。\\[6pt]联立后令 $\Delta_y=0$，\\得到两个参数 $m_1,m_2$。\\[6pt]重根 $y=pm$，切点为\\$Q(pm^2/2,pm)$。\\[6pt]本例 $m=\pm1$，得到\\$Q(1,\pm2)$，连成红线。};
\node[anchor=west,blue] at (-1.6,-3.25) {$y^2=4x$（精确示例，$p=2$）};
\end{tikzpicture}\end{center}\par\smallskip
\end{minipage}\par

此法不预先调用\textbf{切点弦}结论，也不需要导数。
\begin{description}[style=nextline,leftmargin=2.8em,labelsep=0.5em,itemsep=0.5em]
\item[\textbf{第1步：设经过 $P$ 的切线。}]
把待求切线写成
\[\ell:\ x=x_0+m(y-y_0).\]
\textbf{理由：} 这样设可包含 $x=\text{常数}$ 的竖直切线，尤其不会漏掉顶点处的切线 $x=0$。对 $y^2=2px$，切线不会是水平直线。
\item[\textbf{第2步：联立抛物线。}] 把 $\ell$ 代入 $y^2=2px$：
\[
 y^2-2pmy-2p(x_0-my_0)=0.
\]
\item[\textbf{第3步：利用相切条件。}]
直线与抛物线相切时，联立所得关于 $y$ 的一元二次方程有两个相等的实根，因此判别式为零：
\[
 \Delta=(-2pm)^2+8p(x_0-my_0)=0
 \iff\boxed{pm^2-2y_0m+2x_0=0}.
\]
这也是关于 $m$ 的二次方程。其判别式
\[\Delta_m=4y_0^2-8px_0=4(y_0^2-2px_0)>0,\]
故存在两个不同实数 $m_1,m_2$，对应两条不同切线。
\item[\textbf{第4步：找出切点坐标。}]
判别式为零时，$y$ 的二次方程有重根 $y=pm$；再代入抛物线，切点为
\[Q=\left(\frac{pm^2}{2},\,pm\right).\]
\item[\textbf{第5步：证明每个切点落在同一条直线上。}]
由第3步的 $pm^2-2y_0m+2x_0=0$ 得
\[
 y_0(pm)=p\left(x_0+\frac{pm^2}{2}\right).
\]
即切点 $Q$ 满足 $y_0y=p(x+x_0)$。两个不同的 $m$ 分别产生不同切点，故两切点所在直线就是
\[\boxed{y_0y=p(x+x_0)}.\]
\end{description}
\textbf{方法价值：} 遇到不方便直接套切点弦公式的题，也可以依靠\textbf{直线设法 $\to$ 联立 $\to$ 判别式为零 $\to$ 回代切点}解决。

\medskip
\par\noindent\begin{minipage}{\linewidth}
\textbf{\textcolor{CProof}{推广：$x^2=2py$ 的情况}}
\begin{center}\begin{tikzpicture}[x=1cm,y=1cm,>=Stealth,font=\small]
\path[use as bounding box] (-3.5,-2.2) rectangle (9,3.6);
\node[anchor=west,font=\small\bfseries] at (-3.3,3.2) {交换坐标：上开口抛物线的切点弦为水平线};
\draw[->,gray] (-3.1,0)--(3.1,0) node[right] {$x$};
\draw[->,gray] (0,-1.8)--(0,2.4) node[right] {$y$};
\draw[blue,thick,domain=-3:3,samples=81] plot (\x,{\x*\x/4});
\draw[orange!85!black,thick] (-2.7,1.7)--(0,-1)--(2.7,1.7);
\draw[red!75!black,very thick] (-2.8,1)--(2.8,1);
\fill (0,-1) circle(2pt) node[below right] {$P(0,-1)$};
\fill (-2,1) circle(2pt) node[below left,yshift=-4pt] {$A(-2,1)$};
\fill (2,1) circle(2pt) node[below right,yshift=-4pt] {$B(2,1)$};
\node[anchor=north west,align=left,text width=4.4cm] at (4,2.4) {$x^2=4y$，$p=2$。\\[7pt]切线：$y=-x-1$、\\$y=x-1$，均经过 $P$。\\[7pt]切点弦：$y=1$。\\[7pt]一般形式：\\$x_0x=p(y+y_0)$。};
\end{tikzpicture}\end{center}\end{minipage}\par
交换上述各式中的 $x,y$ 即可。已知 $P(x_0,y_0)$ 且 $x_0^2-2py_0>0$，切点处切线为
\[u x=p(y+v),\qquad Q(u,v)\in C,\]
因为经过 $P$，所以 $u x_0=p(y_0+v)$。两个切点都在同一条直线
\[\boxed{x_0x=p(y+y_0)}\]
上，故它就是该抛物线的切点弦所在直线。

\medskip
\textbf{\textcolor{CProof}{条件为什么恰好是大于零（快速自查）}}
设切点 $Q(t^2/(2p),t)$，由该点处切线经过 $P$ 得
\[
 t^2-2y_0t+2px_0=0,\qquad
 \Delta=4(y_0^2-2px_0).
\]
大于零才有两个不同的实数 $t$（也就是两个不同切点）；等于零只有一个切点，小于零没有实切点。纵轴形式同理。
\end{proofbox}
